% 第6章第2节：基于可逆矩阵的对角化（特征基下的矩阵与幂、可对角化的定义与判据及不能对角化的情形、步骤与算例）。

\section{对角化}
\label{sec:ls-matrix-powers}
\label{sec:ch06-powers-spectrum}

上一节找到了被算子只伸缩的方向，并证明属于不同特征值的特征向量线性无关。本节回答由此引出的问题：有了由特征向量组成的基，算子的矩阵变成什么样，又能怎样利用？这样的基什么时候存在？换到这组基所用的换基矩阵只要求可逆，不要求各列互相垂直，所以本节的对角化是用一般的可逆矩阵作相似变换。

\subsection{特征向量基下的矩阵表示}
\label{sec:matrix-analysis-spectrum}

设$\bm v_1,\ldots,\bm v_n$是由算子$T$的特征向量组成的基，$T\bm v_j=\lambda_j\bm v_j$。按第~\ref{sec:ls-map-representation}节的约定，$T$在这组基下的矩阵的第$j$列是$T\bm v_j=\lambda_j\bm v_j$的坐标，即第$j$个位置为$\lambda_j$、其余位置为零。所以这个矩阵在主对角线以外全为零，每个基方向被独立地伸缩。

示例~\ref{ex:ls-operator-eigenvector-check}中的$(\bm e_1,(1,1)^{\top})$就是这样一组基。换到这组基，正是第~\ref{sec:ls-similarity}节的相似变换$\bm P^{-1}\bm A\bm P=\operatorname{diag}(2,3)$；在图~\ref{fig:ls-similarity-change-of-basis-flow}中，基向量围成的平行四边形只沿两条特征方向分别伸长$2$倍和$3$倍。

这一观察也可以写成矩阵等式。

设$\bm A$是$T$在标准基下的矩阵，把特征向量的标准坐标排成$\bm V=[\bm v_1,\ldots,\bm v_n]$。$\bm V$的各列是新基向量的旧坐标，所以它正是第~\ref{sec:ls-similarity}节的换基矩阵；各列线性无关，所以$\bm V$可逆（第~\ref{sec:ls-injective-surjective-invertible}节）。按乘积的列读法，$\bm A\bm V$的第$j$列是$\bm A\bm v_j=\lambda_j\bm v_j$；右乘对角矩阵$\bm\Lambda=\operatorname{diag}(\lambda_1,\ldots,\lambda_n)$把第$j$列乘以$\lambda_j$，所以$\bm V\bm\Lambda$的第$j$列也是$\lambda_j\bm v_j$。于是$\bm A\bm V=\bm V\bm\Lambda$，两边左乘或右乘$\bm V^{-1}$，即得
\[
 \begin{gathered}
\bm V^{-1}\bm A\bm V=\bm\Lambda\\
\bm A=\bm V\bm\Lambda\bm V^{-1}.
 \end{gathered}
\]

对角形式的好处在反复作用时最明显。同一个方阵反复作用，是矩阵乘法最常见的用法之一：定常迭代$\bm x^{(k+1)}=\bm A\bm x^{(k)}$走$k$步后得到$\bm x^{(k)}=\bm A^k\bm x^{(0)}$。第~\ref{sec:ls-matrix-power-rules}节规定了矩阵幂的写法，但直接连乘会不断混合各个坐标，看不出$\bm A^k$在$k$很大时是什么样子。把$\bm A=\bm V\bm\Lambda\bm V^{-1}$连乘$k$次，相邻的$\bm V^{-1}\bm V$相消（这就是式~\eqref{eq:ls-similarity-product}后面的$\bm B^k=\bm P^{-1}\bm A^k\bm P$），得到
\begin{equation}
\begin{gathered}
\bm A^k=\bm V\bm\Lambda^k\bm V^{-1}\\
\bm\Lambda^k=\operatorname{diag}(\lambda_1^k,\ldots,\lambda_n^k).
\end{gathered}
\label{eq:ls-diagonalization-power}
\end{equation}
于是求幂只需对对角元逐个乘方。

三个因子可以从右向左读：$\bm V^{-1}\bm x^{(0)}$给出初值在各特征方向上的坐标，$\bm\Lambda^k$把第$j$个坐标乘以$\lambda_j^k$，$\bm V$再把结果合回标准坐标。上一节已经看到，沿特征方向有$\bm A^k\bm v=\lambda^k\bm v$；式~\eqref{eq:ls-diagonalization-power}把这一点推广到任意初值。

下面用一个上三角矩阵演示这种计算：先把一个初值拆成特征方向上的分量，再对所有初值一并写出矩阵幂。
\begin{example}[用特征向量计算矩阵幂]
\label{ex:ls-triangular-power}
给定
\[
 \bm A=\begin{bmatrix}0.8&4\\0&0.5\end{bmatrix},
\]
计算$\bm A^k$。

第一步，求特征值与特征向量。$\bm A$是上三角矩阵，特征值为对角元$0.8$与$0.5$。直接验证
\[
\begin{gathered}
\bm A\begin{bmatrix}1\\0\end{bmatrix}=0.8\begin{bmatrix}1\\0\end{bmatrix}\\
\bm A\begin{bmatrix}-40/3\\1\end{bmatrix}
 =\begin{bmatrix}-32/3+4\\0.5\end{bmatrix}
 =0.5\begin{bmatrix}-40/3\\1\end{bmatrix},
\end{gathered}
\]
所以$\bm v_1=(1,0)^{\top}$与$\bm v_2=(-40/3,1)^{\top}$是特征向量。

第二步，计算一条轨迹。初值$(0,1)^{\top}$不在特征方向上，先把它拆成$(0,1)^{\top}=\tfrac{40}{3}\bm v_1+\bm v_2$。两个分量每作用一次分别乘以$0.8$与$0.5$，由线性性，
\[
 \bm A^k\begin{bmatrix}0\\1\end{bmatrix}
 =\frac{40}{3}\,0.8^k\bm v_1+0.5^k\bm v_2
 =\begin{bmatrix}\tfrac{40}{3}(0.8^k-0.5^k)\\0.5^k\end{bmatrix}.
\]
$k=1$时得$(4,0.5)^{\top}$，与直接计算$\bm A(0,1)^{\top}$一致。

第三步，对所有初值一并计算，即应用式~\eqref{eq:ls-diagonalization-power}。取$\bm V=\left[\begin{smallmatrix}1&-40/3\\0&1\end{smallmatrix}\right]$，则$\bm V^{-1}=\left[\begin{smallmatrix}1&40/3\\0&1\end{smallmatrix}\right]$，
\[
 \bm A^k=\bm V\begin{bmatrix}0.8^k&0\\0&0.5^k\end{bmatrix}\bm V^{-1}
 =\begin{bmatrix}0.8^k&\tfrac{40}{3}(0.8^k-0.5^k)\\0&0.5^k\end{bmatrix}.
\]
结论：$\bm A^k$的第二列正是第二步的轨迹；整个迭代拆成两个独立的数列，分别按$0.8^k$与$0.5^k$衰减。
\end{example}
这种计算依赖一个前提：特征向量足以组成整个空间的一组基。下一小节讨论这个前提何时成立。

\subsection{可对角化的条件与边界}

\paragraph{定义与特征向量判据}
上一小节说明，有了由特征向量组成的基，$\bm A$就相似于一个对角矩阵。把相似于对角矩阵这一性质取作定义。
\begin{definition}[实可对角化]
\label{def:ls-real-diagonalizable}
设$\bm A$是$n$阶实方阵。若存在$n$阶实可逆矩阵$\bm V$与实对角矩阵$\bm\Lambda$使$\bm V^{-1}\bm A\bm V=\bm\Lambda$，即$\bm A$相似于实对角矩阵，就称$\bm A$\term{实可对角化}（diagonalizable over the reals）；求出这样的$\bm V$与$\bm\Lambda$，称为把$\bm A$对角化。
\end{definition}
“实”表示$\bm V$与$\bm\Lambda$的元素都是实数。本书在实数范围内讨论，下文所说的可对角化都指实可对角化。

哪些矩阵可以对角化？下面的判据把问题化为特征向量够不够多。
\begin{theorem}[对角化判据]
\label{thm:ls-diagonalization-criterion}
$n$阶实方阵$\bm A$实可对角化，当且仅当$\bm A$有$n$个线性无关的实特征向量。这时以这些特征向量为列组成$\bm V$，以相应特征值按同一顺序组成$\bm\Lambda=\operatorname{diag}(\lambda_1,\ldots,\lambda_n)$，就有$\bm V^{-1}\bm A\bm V=\bm\Lambda$。
\end{theorem}
\begin{proof}
设$\bm V$的各列为$\bm v_1,\ldots,\bm v_n$，$\bm\Lambda=\operatorname{diag}(\lambda_1,\ldots,\lambda_n)$。上一小节的列读法对任意这样的$\bm V$与$\bm\Lambda$都成立，它给出
\[
 \begin{gathered}
\bm A\bm V=\bm V\bm\Lambda
 \quad\Longleftrightarrow\quad
 \bm A\bm v_j=\lambda_j\bm v_j\\
j=1,\ldots,n.
 \end{gathered}
\]
若$\bm A$有$n$个线性无关的特征向量，以它们为列的$\bm V$可逆，左乘$\bm V^{-1}$即得$\bm V^{-1}\bm A\bm V=\bm\Lambda$。反过来，若$\bm V^{-1}\bm A\bm V=\bm\Lambda$，左乘$\bm V$得$\bm A\bm V=\bm V\bm\Lambda$，于是每列都满足$\bm A\bm v_j=\lambda_j\bm v_j$；$\bm V$可逆，所以各列线性无关，从而都非零，它们是$n$个线性无关的特征向量。
\end{proof}
由第~\ref{sec:ls-eigenvector-independence}节，属于不同特征值的特征向量线性无关，由此得到一个常用的充分条件：$n$阶实方阵若有$n$个互不相同的实特征值，就实可对角化。示例~\ref{ex:ls-triangular-power}的矩阵即属此类。这个条件不是必要的：单位矩阵$\bm I_n$只有一个特征值$1$，本身却已经是对角矩阵。

\paragraph{不能对角化的矩阵}
判据要求$n$个线性无关的实特征向量，这一要求可能在两处得不到满足。其一，实特征值不够。示例~\ref{ex:ls-rotation-no-real-eigenvector}与示例~\ref{ex:ls-characteristic-polynomial-two-by-two}已经说明，平面上的$90^\circ$旋转$\bm R=\left[\begin{smallmatrix}0&-1\\1&0\end{smallmatrix}\right]$没有实特征值，也就没有实特征向量，所以不能实对角化。

其二，实特征值有了，特征向量却不够。示例~\ref{ex:ls-multiplicity-gap}中的$\bm J=\left[\begin{smallmatrix}0.8&1\\0&0.8\end{smallmatrix}\right]$只有特征值$0.8$，而特征向量只有$(1,0)^{\top}$的倍数，凑不出两个线性无关的特征向量，所以$\bm J$也不能对角化。

对这样的矩阵，式~\eqref{eq:ls-diagonalization-power}不再可用，幂也不再是纯粹的乘方。下面直接计算$\bm J^k$。
\begin{example}[不能对角化矩阵的幂]
\label{ex:ls-jordan-block-power}
给定$\bm J=\left[\begin{smallmatrix}0.8&1\\0&0.8\end{smallmatrix}\right]$，计算$\bm J^k$。把$\bm J$写成$0.8\bm I+\bm N$，其中
\[
\begin{gathered}
\bm N=\begin{bmatrix}0&1\\0&0\end{bmatrix}\\
\bm N^2=\bm0 .
\end{gathered}
\]
$0.8\bm I$与$\bm N$可交换，所以二项式展开成立；又因为$j\geq2$时$\bm N^j=\bm0$，展开后只剩前两项：
\[
 \bm J^k=(0.8\bm I+\bm N)^k=0.8^k\bm I+k\,0.8^{k-1}\bm N=\begin{bmatrix}0.8^k&k\,0.8^{k-1}\\0&0.8^k\end{bmatrix}.
\]
右上角相邻两项之比为$\dfrac{(k+1)\,0.8^{k}}{k\,0.8^{k-1}}=0.8\cdot\dfrac{k+1}{k}$，它在$k<4$时大于$1$，在$k>4$时小于$1$。结论：右上角的$k\,0.8^{k-1}$随$k$先增大、后衰减，不能由单个特征值的乘方描述。
\end{example}

\paragraph{重数判据}
两个反例分别缺了什么？$\bm R$的特征多项式$z^2+1$不能分解为实一次因子；$\bm J$的特征多项式$(z-0.8)^2$分解了，但特征值$0.8$的几何重数$1$小于代数重数$2$。下面说明，这两种缺陷恰好是全部障碍。于是能否对角化可以直接由特征多项式和特征空间的维数判断，而不必逐个尝试特征向量。

先证明几何重数不超过代数重数。

设特征值$\lambda$的特征空间维数为$g$，取它的一组基，并扩充成全空间的一组基。在新基下，前$g$个基向量都被送到自身的$\lambda$倍，所以$\bm A$相似于分块上三角矩阵$\left[\begin{smallmatrix}\lambda\bm I_g&*\\\bm0&\bm C\end{smallmatrix}\right]$。相似不改变特征多项式（第~\ref{sec:ls-eigen-similarity-invariance}节），所以可以用这个分块矩阵计算$p_{\bm A}$。用$z\bm I$减去它以后，前$g$列都只在对角位置上有非零元素$z-\lambda$。沿第一列按余子式展开（第~\ref{sec:ls-determinant-properties}节），得到$z-\lambda$乘以一个形状相同、但只剩$g-1$个这种列的行列式；重复$g$次，即得$p_{\bm A}(z)=(z-\lambda)^g\,p_{\bm C}(z)$。所以代数重数至少为$g$。示例~\ref{ex:ls-multiplicity-gap}说明不等号可以严格成立。

由此得到可以计算的判据：$\bm A$实可对角化，当且仅当$p_{\bm A}$能分解为实一次因子之积，且每个特征值的几何重数等于代数重数。

先证充分性。$p_{\bm A}$分解为$n$个实一次因子时，各特征值的代数重数之和为$n$；几何重数与代数重数逐个相等，所以几何重数之和也是$n$。由第~\ref{sec:ls-eigenvector-independence}节，不同特征值的特征空间之和是直和，各特征空间的基合在一起仍线性无关，于是恰好得到$n$个线性无关的特征向量；由对角化判据，$\bm A$可对角化。

再证必要性。若$\bm A=\bm V\bm\Lambda\bm V^{-1}$，则由相似不变性，$p_{\bm A}=p_{\bm\Lambda}=\prod_i(z-\lambda_i)$已分解为实一次因子。对每个特征值$\lambda$，它在$\bm\Lambda$对角线上出现的次数既是代数重数，也是$\bm\Lambda-\lambda\bm I$的零空间维数，即$\bm\Lambda$的几何重数。由第~\ref{sec:ls-eigen-similarity-invariance}节，相似不改变几何重数，所以$\bm A$的两种重数也相等。

有$n$个互不相同的实特征值时，每个代数重数都是$1$，几何重数至少为$1$又不超过$1$，判据的条件自动满足，这就回到了前面的充分条件。

\subsection{对角化的计算与用途}

按判据，对角化分四步。第一步，求出特征多项式$p_{\bm A}$的全部实根，即全部实特征值。第二步，对每个特征值$\lambda$解齐次方程$(\bm A-\lambda\bm I)\bm v=\bm0$，得到特征空间的一组基。第三步，若这些向量合起来共有$n$个，就以它们为列组成$\bm V$，以对应的特征值按同一顺序组成$\bm\Lambda$。第四步，核对$\bm A\bm V=\bm V\bm\Lambda$。若特征多项式不能分解为实一次因子，或某个特征空间的维数小于相应的代数重数，过程就在第一步或第三步停止，矩阵不能实对角化，这正是上一小节的两种情形。

下面用一个两状态转移的例子走一遍这四步。

设一个群体在两个状态之间转移：每期处于状态一的个体有$10\%$转入状态二，处于状态二的个体有$20\%$转入状态一。记第$k$期两个状态所占的比例为$\bm x^{(k)}=(x_1^{(k)},x_2^{(k)})^{\top}$，则下一期状态一的比例为$0.9x_1^{(k)}+0.2x_2^{(k)}$，状态二的比例为$0.1x_1^{(k)}+0.8x_2^{(k)}$，即
\[
\begin{gathered}
\bm x^{(k+1)}=\bm A\bm x^{(k)}\\
\bm A=\begin{bmatrix}0.9&0.2\\0.1&0.8\end{bmatrix}.
\end{gathered}
\]
\begin{example}[两状态转移矩阵的对角化]
\label{ex:ls-transition-matrix-diagonalization}
给定转移矩阵$\bm A=\left[\begin{smallmatrix}0.9&0.2\\0.1&0.8\end{smallmatrix}\right]$，按四步把它对角化。

第一步，求特征值：
\[
 p_{\bm A}(z)=(z-0.9)(z-0.8)-0.2\cdot0.1=z^2-1.7z+0.7=(z-1)(z-0.7),
\]
得$\lambda_1=1$、$\lambda_2=0.7$。两个特征值不同，所以$\bm A$可对角化。

第二步，求特征空间。$\lambda_1=1$时$\bm A-\bm I=\left[\begin{smallmatrix}-0.1&0.2\\0.1&-0.2\end{smallmatrix}\right]$，两行成比例，方程化为$-0.1x_1+0.2x_2=0$，取$\bm v_1=(2,1)^{\top}$。$\lambda_2=0.7$时$\bm A-0.7\bm I=\left[\begin{smallmatrix}0.2&0.2\\0.1&0.1\end{smallmatrix}\right]$，方程化为$x_1+x_2=0$，取$\bm v_2=(1,-1)^{\top}$。

第三步，组成
\[
\begin{gathered}
\bm V=\begin{bmatrix}2&1\\1&-1\end{bmatrix}\\
\bm\Lambda=\begin{bmatrix}1&0\\0&0.7\end{bmatrix}\\
\bm V^{-1}=\frac13\begin{bmatrix}1&1\\1&-2\end{bmatrix},
\end{gathered}
\]
其中$\bm V^{-1}$由二阶逆矩阵公式得到，$\det\bm V=-3$。

第四步，核对：
\[
\begin{gathered}
\bm A\bm v_1=\begin{bmatrix}1.8+0.2\\0.2+0.8\end{bmatrix}=\bm v_1\\
\bm A\bm v_2=\begin{bmatrix}0.9-0.2\\0.1-0.8\end{bmatrix}=0.7\bm v_2,
\end{gathered}
\]
所以$\bm A\bm V=\bm V\bm\Lambda$。也可以用上一节的系数关系作快速检查：$\operatorname{tr}\bm A=1.7=1+0.7$，$\det\bm A=0.7=1\cdot0.7$。

结论：$\bm A=\bm V\bm\Lambda\bm V^{-1}$，$\bm V$的两列分别是属于特征值$1$与$0.7$的特征向量。
\end{example}

有了这个分解，就可以按式~\eqref{eq:ls-diagonalization-power}回答这个群体的长期去向。
\begin{example}[转移矩阵的长期比例]
\label{ex:ls-transition-matrix-long-run}
沿用示例~\ref{ex:ls-transition-matrix-diagonalization}的分解，计算$\bm A^k$以及$k\to\infty$时两个状态的比例。由式~\eqref{eq:ls-diagonalization-power}，
\[
 \bm A^k=\bm V\begin{bmatrix}1&0\\0&0.7^k\end{bmatrix}\bm V^{-1}
 =\frac13\begin{bmatrix}2+0.7^k&2-2\cdot0.7^k\\1-0.7^k&1+2\cdot0.7^k\end{bmatrix};
\]
$k=1$时右端正是$\bm A$。若开始时全部个体处于状态一，即$\bm x^{(0)}=(1,0)^{\top}$，则$\bm x^{(k)}$是$\bm A^k$的第一列：
\[
 \bm x^{(k)}=\frac13\begin{bmatrix}2+0.7^k\\1-0.7^k\end{bmatrix}
 \to\begin{bmatrix}2/3\\1/3\end{bmatrix}\qquad(k\to\infty).
\]
对一般的初始比例$\bm x^{(0)}=(a,1-a)^{\top}$，$\bm V^{-1}\bm x^{(0)}$给出它在两条特征方向上的分量：
\[
 \bm V^{-1}\bm x^{(0)}=\frac13\begin{bmatrix}a+(1-a)\\a-2(1-a)\end{bmatrix}=\begin{bmatrix}1/3\\(3a-2)/3\end{bmatrix}.
\]
沿$\bm v_1$的分量总是$1/3$，沿$\bm v_2$的分量每期乘以$0.7$而趋于零。结论：无论从什么比例出发，长期比例都是$\tfrac13\bm v_1=(2/3,1/3)^{\top}$。
\end{example}

至此，本节开头的问题有了回答。在由特征向量组成的基下，算子的矩阵是对角的；以特征向量为列的$\bm V$把$\bm A$化为$\bm\Lambda$，反复作用化为特征值的乘方，$\bm A^k=\bm V\bm\Lambda^k\bm V^{-1}$。$n$阶矩阵可对角化，当且仅当它有$n$个线性无关的实特征向量；等价地，特征多项式分解为实一次因子，且每个特征值的几何重数等于代数重数。

实特征值不足或特征向量不足时，矩阵不能对角化；后一种情形的幂含有随$k$增长的多项式因子。

对角化回答了“能否找到只被伸缩的基方向”，却没有要求这些方向互相垂直。示例~\ref{ex:ls-transition-matrix-diagonalization}中的两条特征方向$(2,1)^{\top}$与$(1,-1)^{\top}$内积为$1$，并不垂直。换到这样的基，求新坐标要解方程，新坐标的平方和也不等于原向量长度的平方。下一节追问，哪些矩阵能用正交矩阵对角化，即具有一组标准正交的特征向量。至于不能对角化的矩阵在标准正交基下能化到多简单，以及特征向量接近线性相关时矩阵幂有多大，分别留给第~\ref{sec:ls-general-square-spectrum}节与第~\ref{sec:ls-matrix-power-behavior}节。
